% ============================================================
% SM-10: First-Principles Quark Mass from Finite Element
% Chain Network Simulation
% Version 0.1 — 9 April 2026 (PROPOSAL)
%
% CHANGELOG:
%   v0.1  9 Apr 2026 — Proposal draft (Claude Opus)
%
% CONTRIBUTORS:
%   Thomas Lee Abshier ND — Physical vision: chain-type
%     decomposition, pine tree model, three-region bonding,
%     surface blanket, FEM proposal
%   Claude Opus (Anthropic) — Chain-type computation,
%     cooperative enhancement analysis, proposal drafting
% ============================================================

\documentclass[11pt,letterpaper]{article}
\usepackage[utf8]{inputenc}
\usepackage[margin=1in]{geometry}
\usepackage[T1]{fontenc}
\usepackage{amsmath,amssymb,amsfonts,amsthm}
\usepackage{graphicx}
\usepackage{xcolor}
\usepackage{hyperref}
\usepackage{booktabs}
\usepackage{array}
\usepackage{float}
\usepackage[font=small, labelfont=bf]{caption}
\usepackage{parskip}
\usepackage[authoryear,round]{natbib}

\hypersetup{
    colorlinks=true,
    linkcolor=blue,
    citecolor=blue,
    urlcolor=blue
}

\newtheorem{theorem}{Theorem}[section]
\newtheorem{remark}[theorem]{Remark}

\newcommand{\phig}{\varphi}

\title{\textbf{SM-10: First-Principles Quark Mass from\\Finite Element Chain Network Simulation}\\[6pt]
{\large 600-Cell Standard Model Emergence Series}\\[4pt]
{\normalsize Version 0.1 (Proposal), 9 April 2026}}

\author{Thomas Lee Abshier, ND \and Claude Opus (Anthropic)}

\date{\normalsize Hyperphysics Institute\\
\url{https://hyperphysics.com}\\
\href{mailto:drthomas007@protonmail.com}{drthomas007@protonmail.com}}

\begin{document}
\maketitle

\begin{abstract}
SM-8 established a zero-free-parameter quark mass formula
$M_q = m_e(z/\phig)V^{7/3}$ that predicts all four heavy quark
masses to RMS~$2.1\%$.  SM-9 showed that the exponent~$7/3$
arises from pair counting $\times$ cage linear dimension, and
that the mass is physically distributed across a network of
radial and tangential DP chains filling the cage interior.
However, the exponent has not been derived from first principles.
This paper proposes a finite element method (FEM) simulation to
compute quark masses directly: place cage CPs at their
lattice-determined positions, fill the interior with DP Sea,
let every CP seek opposite-polarity targets to form chains, and
sum the total organised DP count~$\times$~$M_0$.  If the
simulation reproduces $V^{7/3}$ scaling without imposing any
power law, this would constitute the first derivation of quark
mass ratios from geometry alone.
\end{abstract}

\tableofcontents

% ============================================================
\section{Motivation}
% ============================================================

The CPP quark mass programme has progressed through three stages:
\begin{enumerate}
\item \textbf{SM-8 v2.1:} Calibrated formula
  $m \sim V^{2.38} \times z$ (2~parameters, 0.02\% top precision).
\item \textbf{SM-8 v4.0:} Zero-parameter formula
  $M = m_e(z/\phig)V^{7/3} \times [z C_F]$ (0~parameters, 2.1\% RMS).
\item \textbf{SM-9 v2.0:} Partial derivation of $7/3$ from pair
  counting; chain-type energy budget decomposition.
\end{enumerate}

The remaining gap: the exponent~$7/3$ is motivated but not
rigorously derived.  The chain-type analysis (SM-9 v2.0) showed
that the mass arises from a complex volume-filling network of DP
chains with three distinct bonding regions, but the exact scaling
of this network with cage size has not been computed from first
principles.

A direct simulation---placing CPs, forming chains, counting
organised DPs---would bypass the need for an analytical
derivation and provide a numerical first-principles prediction.

% ============================================================
\section{Physical Model}
\label{sec:model}
% ============================================================

\subsection{The chain network}

A caged quark consists of:
\begin{itemize}
\item A central qCP at the cage centre.
\item $V_{\rm opp}$ opposite-polarity cage vertices (chain anchors).
\item $V_{\rm same}$ same-polarity cage vertices (surface radiators).
\item $E_{\rm attr}$ attractive cage edges (tangential chain sites).
\item The DP Sea filling the interior volume.
\end{itemize}

Each CP (whether on a chain or in the Sea) has two charges of
opposite polarity.  Each charge seeks the nearest opposite-polarity
target, forming a chain link.  This process cascades: every CP
recruited into a chain becomes a new source of lateral connections.

\subsection{Three bonding regions}

The chain network self-organises into three regions
(identified by Abshier):

\begin{description}
\item[Region 1 (near central CP):]
  Radial chains converge at small separations.  Tangential CPs
  terminate on the central CP or adjacent radial CPs.  Dense
  cross-linking.

\item[Region 2 (mid-cage):]
  Radials have diverged.  Tangential chains extend to
  neighbouring radials' tangential chains, forming a web mesh.
  Every DP's two CPs seek opposite-polarity targets in all
  directions---the ``complex melee.''

\item[Region 3 (near cage surface):]
  Tangential chains arch upward toward opposite-polarity cage
  surface CPs.  Organisation converges on cage vertices.
\end{description}

\subsection{Surface blanket}

Same-polarity cage vertices ($V_{\rm same}$) launch outward
radial chains to the thermalization distance.  These surface
radials develop their own tangential linking, creating a
structure identical to the up quark's radial-plus-tangential
blanket.  This contributes additional organised DPs.

\subsection{Linear oscillator (down-type quarks)}

Down-type quarks ($d$, $s$, $b$) include a captured eCP that
adds a linear/radial ZBW oscillation mode.  This mode contributes
mass through its own chain structure.

\subsection{Mass calculation}

The quark mass equals the total number of organised DPs in the
chain network multiplied by the bare DP energy:
\begin{equation}
M_q = M_0 \times N_{\rm organised}
\end{equation}
where $M_0 = m_e \times z/\phig = 3.790$~MeV is the energy per
DP link and $N_{\rm organised}$ is the total count of DPs
recruited from the Sea into the chain network.

% ============================================================
\section{Simulation Design}
\label{sec:simulation}
% ============================================================

\subsection{Geometry setup}

\begin{enumerate}
\item Place a central qCP at the origin.
\item Place cage CPs at the lattice-determined positions for the
  target cage (tetrahedron, icosahedron, dodecahedron, or
  icosidodecahedron).
\item Assign charge polarity: $V/2$ positive, $V/2$ negative,
  at the $2/3$ attractive fraction assignment.
\item Fill the interior volume with DP Sea at natural density
  $\rho_{\rm Sea}$ (to be determined from the lattice spacing
  $l_P$ or parameterised).
\end{enumerate}

\subsection{Chain formation algorithm}

At each simulation step:
\begin{enumerate}
\item For each CP (central, cage, and Sea), identify the nearest
  opposite-polarity CP that is not already bonded.
\item Form a chain link (DP bond) between them.
\item Mark both CPs as ``organised'' (part of the chain network).
\item Newly organised CPs become sources for the next step's
  lateral connections.
\end{enumerate}

The simulation terminates when no further chain formation occurs
(all nearby opposite-polarity targets are bonded or beyond the
thermalization distance).

\subsection{Boundary conditions}

\begin{itemize}
\item \textbf{Inner boundary:} The central CP is fixed.
\item \textbf{Cage boundary:} Cage CPs are fixed at their
  lattice positions.  They act as impedance boundaries
  (Appendix~B of SM-8).
\item \textbf{Outer boundary:} DPs beyond the thermalization
  distance $r_{\rm therm}$ are unorganised (remain part of the
  random Sea).  For the top quark, the Shell~3 gap is
  automatically present in the lattice geometry.
\end{itemize}

\subsection{Output}

The primary output is $N_{\rm organised}$: the total count of DPs
recruited into the chain network.  Secondary outputs include:
\begin{itemize}
\item Spatial distribution of organised DPs (Region 1/2/3 fractions)
\item Radial vs.\ tangential chain ratio
\item Surface blanket DP count
\item Chain length distribution
\item Effective fractal dimension of the network
\end{itemize}

% ============================================================
\section{Predictions and Success Criteria}
\label{sec:predictions}
% ============================================================

\subsection{Primary prediction}

If the simulation correctly models the physics, the organised DP
counts should satisfy:
\begin{equation}
\frac{N_{\rm organised}(q_1)}{N_{\rm organised}(q_2)}
= \frac{M_{q_1}}{M_{q_2}}
= \frac{V_1^{7/3} \times {\rm gap}_1}{V_2^{7/3} \times {\rm gap}_2}
\end{equation}
for all quark pairs.  Specifically:

\begin{table}[H]
\centering
\caption{Target DP count ratios (normalised to strange).}
\begin{tabular}{@{}lrrr@{}}
\toprule
Ratio & Target & From $V^{7/3}$ & From PDG \\
\midrule
$N_c / N_s$ & 13.0 & 13.0 & 13.6 \\
$N_b / N_s$ & 42.8 & 42.8 & 44.8 \\
$N_t / N_s$ & 1762 & 1762 & 1850 \\
\bottomrule
\end{tabular}
\end{table}

\subsection{Secondary predictions}

\begin{enumerate}
\item The radial/tangential energy split should be $\sim$44\%/56\%
  across all cages (from SM-9 chain-type analysis).
\item The three-region structure should emerge naturally without
  being imposed.
\item The effective fractal dimension of the chain network should
  be between 2 and 3 (corresponding to the $7/3$ exponent).
\item The top quark simulation should show qualitatively different
  chain structure (long-range coordination bonds through the
  Shell~3 gap).
\end{enumerate}

\subsection{Success criteria}

\begin{description}
\item[Level A (strong success):] DP count ratios match PDG mass
  ratios to $< 5\%$.  The $V^{7/3}$ scaling emerges without being
  imposed.  Fractal dimension measured.
\item[Level B (partial success):] DP count ratios show the correct
  ordering and approximate magnitude.  The three-region structure
  emerges.  Scaling exponent is measurable but differs from $7/3$.
\item[Level C (informative failure):] The simulation reveals which
  physical assumptions of the chain model are incorrect, guiding
  revision of SM-9.
\end{description}

% ============================================================
\section{Computational Requirements}
\label{sec:computation}
% ============================================================

\subsection{Scale estimates}

\begin{itemize}
\item \textbf{Strange quark (tetrahedron):}
  Confinement sphere radius $\sim 1/\phig$.
  Volume $\sim 1$ (lattice units$^3$).
  At $\rho_{\rm Sea} \sim 100$~DPs per unit volume:
  $\sim 100$ DPs.  Trivial on CPU.

\item \textbf{Top quark (icosidodecahedron):}
  Confinement sphere radius $\sim \sqrt{2}$.
  Volume $\sim 12$ (lattice units$^3$).
  At same density: $\sim 1200$ DPs.
  Plus surface blanket: $\sim 2000$ total.  Still CPU-tractable.

\item \textbf{With fine-grained DP Sea} ($\rho \sim 10^4$):
  $\sim 10^4$--$10^5$ DPs.  GPU recommended.
\end{itemize}

\subsection{Implementation options}

\begin{description}
\item[Phase 1 (CPU, proof of concept):]
  Coarse-grained simulation with $\rho \sim 100$.
  Test whether DP count ratios track $V^{7/3}$.
  Python with NumPy.  Estimated time: 1--2 sessions.

\item[Phase 2 (GPU, production):]
  Fine-grained simulation with $\rho \sim 10^4$.
  Measure fractal dimension, region fractions, chain statistics.
  CUDA or JAX.  Estimated time: 1--2 weeks.

\item[Phase 3 (lattice-scale):]
  Full 600-cell tessellation with explicit lattice CPs.
  Test whether the simulation reproduces the Shell~3 gap effect
  and the $z \times C_F = 16$ multiplier.
  GPU cluster.  Estimated time: months.
\end{description}

% ============================================================
\section{Relationship to Existing Methods}
\label{sec:comparison}
% ============================================================

\begin{description}
\item[Lattice QCD:]
  Computes quark masses on a discretised spacetime lattice using
  Monte Carlo sampling of the path integral.  Requires the QCD
  Lagrangian as input.  CPP's FEM simulation uses a different
  dynamical rule (CP charge-seeking rather than gauge field
  evolution) and a different lattice (600-cell rather than
  hypercubic).

\item[Bag models:]
  Confine quarks within a spherical boundary with phenomenological
  boundary conditions.  The CPP cage plays a similar role but the
  boundary is a polyhedral lattice substructure, not imposed.

\item[String models:]
  Model confinement via flux tubes between quarks.  CPP's radial
  chains are the microscopic realisation of the flux tube, with
  the chain composed of explicit DP links.
\end{description}

The CPP FEM simulation is distinguished by:
\begin{enumerate}
\item The lattice is derived from geometry (600-cell), not chosen
  for computational convenience.
\item The dynamical rule (charge-seeking) is simpler than gauge
  field evolution.
\item The mass prediction is a DP count, not an eigenvalue of a
  transfer matrix.
\item Zero free parameters in the mass formula (if the simulation
  reproduces $V^{7/3}$).
\end{enumerate}

% ============================================================
\section{Open Questions for the Simulation}
\label{sec:questions}
% ============================================================

\begin{enumerate}
\item \textbf{What is $\rho_{\rm Sea}$?}
  The DP Sea density should be derivable from the 600-cell lattice
  spacing and the ZBW frequency.  Is it $1/l_{\rm edge}^3$?
  $z/l_{\rm edge}^3$?  This is the only parameter not yet fixed.

\item \textbf{Does the gap multiplier emerge?}
  For the top quark, the Shell~3 gap should force chains to use
  coordination bonds.  Does the simulation naturally produce a
  $\times 16$ enhancement, or does it require explicit modelling
  of the colour factor $C_F$?

\item \textbf{What sets the thermalization distance?}
  Chains terminate when the organised DP's SSV field falls below
  the Sea's thermal noise.  This distance should be computable
  from $\alpha_{\rm geom} = 1/\sqrt{5}$.

\item \textbf{Does the exponent drift?}
  SM-9 documented a 2\% drift in the pairwise exponent across
  pre-gap quarks.  Does the simulation reproduce this drift,
  and does it reveal the physical mechanism?

\item \textbf{Can the simulation distinguish $7/3$ from $3 - 1/\phig$?}
  Grok's conjectured exponent ($2.382$) differs from $7/3$
  ($2.333$) by 2\%.  A simulation with sufficient precision could
  distinguish them.
\end{enumerate}

% ============================================================
\section{Timeline and Milestones}
\label{sec:timeline}
% ============================================================

\begin{table}[H]
\centering
\caption{Proposed milestones for SM-10.}
\begin{tabular}{@{}llll@{}}
\toprule
Phase & Milestone & Deliverable & Est.\ time \\
\midrule
0 & Simulation design document & This paper (v0.1) & Done \\
1a & CPU proof-of-concept & DP counts for all 4 cages & 1--2 sessions \\
1b & Ratio test & Check $N_q$ ratios vs $V^{7/3}$ & Same session \\
2a & GPU implementation & CUDA/JAX chain-formation kernel & 1 week \\
2b & Production runs & 4 cages $\times$ 10 densities & 1 week \\
2c & Analysis & Fractal dimension, region fractions & 1 session \\
3 & Paper draft & SM-10 v1.0 & 1--2 sessions \\
\bottomrule
\end{tabular}
\end{table}

% ============================================================
\section{Conclusion}
\label{sec:conclusion}
% ============================================================

The FEM chain network simulation is the natural next step for the
CPP quark mass programme.  It would convert the scaling formula
$M = m_e(z/\phig)V^{7/3}$ from a calibrated observation into a
first-principles computation, and it would test Thomas Abshier's
physical picture of three-region bonding, fractal cascading, and
surface blanket formation.  If successful, it would be the first
derivation of quark mass ratios from geometry in any theoretical
framework.

The Phase~1 proof-of-concept (CPU, coarse-grained) can be
attempted immediately.  If ratios are promising, Phase~2 (GPU,
fine-grained) would follow.

% ============================================================
\section*{Acknowledgements}
% ============================================================

Thomas Lee Abshier, ND conceived the chain-type decomposition,
the pine tree model with three bonding regions, the surface
blanket structure, and the FEM simulation proposal.

Claude Opus (Anthropic) performed the chain-type energy budget
computation, the cooperative enhancement analysis, the fractal
cascade calculation, and drafted this proposal.

The CPP programme is registered at OSF
(DOI:~\url{https://doi.org/10.17605/OSF.IO/JXE8D}) and maintained
at GitHub (\url{https://github.com/Hyperphysics-Institute/CPP}).

\bibliographystyle{plainnat}
\bibliography{SM-10_references}

\end{document}
